\section{Anticipated results}

\subsection{Phase 2: Sequence processing and demultiplexing}
Successful demultiplexing produces staged Cutadapt inputs, per-step logs and statistics, quality-control summaries, and per-selection count tables that can be used directly for downstream quality control and enrichment analysis:

\begin{lstlisting}
campaign/
|---- config.yaml
|---- demultiplex/
|   |---- cutadapt_input_files/
|   |   |---- 0-S0.fastq
|   |   |---- ...
|   |   |---- 6-S1.fastq
|   |   `---- demultiplex.sh
|   `---- cutadapt_output_files/
|       |---- 0-S0.cutadapt.json
|       |---- 0-S0.cutadapt.log
|       |---- ...
|       |---- input.fastq.gz
|       |---- out.fastq.gz
|       `---- reads_with_adapters.gz
|---- qc/
|   |---- report.txt
|   |---- hits_0-S0.{parquet,pdf,png}
|   |---- ...
|   `---- hits_6-S1.{parquet,pdf,png}
`---- selections/
    |---- AG24_4/counts.txt
    |---- ...
    `---- AG24_4_counts.txt
\end{lstlisting}

For high-quality datasets, the demultiplexing outputs typically show (i) high read retention across the sequential adapter-matching steps, (ii) barcode coverage across most or all expected building blocks, and (iii) low sequencing-error rates in the retained reads. The corresponding quality-control plots (Figures 4 and 5) provide a practical overview of these trends and help distinguish balanced selections from runs affected by barcode-definition errors, poor read quality, or library-construction problems.

\subsection{Phase 4: Selection-focused library reconstruction}
Selection-focused enumeration produces a chemically interpretable view of the enriched subset for a chosen selection together with auxiliary visualization and property files:

\begin{lstlisting}
campaign/
`---- library/
    |---- AG24_4_top_hits.parquet
    |---- library.parquet
    |---- visualization/
    |   |---- reaction_graph.png
    |   |---- reaction_graph_additional.png
    |   |---- reaction_graph_building_blocks.png
    |   |---- compounds/
    |   |   `---- scaffold_1.png
    |   |---- reactions/
    |   |   |---- ABF1.png
    |   |   |---- ABF2.png
    |   |   |---- CuAAC.png
    |   |   `---- SR.png
    |   |---- building_blocks/
    |   |   |---- B0/
    |   |   |   `---- *.png
    |   |   `---- B1/
    |   |       `---- *.png
    |   `---- libraries/
    |       `---- AG24_4_top_hits/
    |           `---- B0=<id>-B1=<id>.png
    `---- properties/
        |---- AG24_4_top_hits/
        |   |---- properties.parquet
        |   `---- prop_*.png
        `---- library/
            |---- properties.parquet
            `---- prop_*.png
\end{lstlisting}

In this phase, the reaction graph (Figure 2) should confirm that the library architecture and reaction order were translated correctly from the workbook definition. The expected outputs are (i) an enumerated catalog linking selection counts to reconstructed structures and (ii) visualization files that allow rapid manual inspection of the top-ranked compounds. For the single-stranded example workflow, the molecular-weight distribution and representative top-hit renderings shown in Figure 3 provide a concise view of the resulting chemical space.

\subsection{Phase 5: Statistical analysis and hit identification}
The example workflow contains outputs for multiple enrichment modes, including count-based comparisons, edgeR-based analysis, and z-score-based ranking:

\begin{lstlisting}
campaign/
|---- analysis/
|   |---- counts/
|   |   `---- condition_vs_control/
|   |       |---- samples.csv
|   |       |---- data.csv
|   |       |---- enrichment_counts.R
|   |       |---- stats.csv
|   |       |---- hits.csv
|   |       |---- condition.csv
|   |       |---- control.csv
|   |       `---- correlation_{condition,control}.png
|   |---- edgeR/
|   |   `---- condition_vs_control/
|   |       |---- samples.csv
|   |       |---- data.csv
|   |       |---- enrichment_edgeR.R
|   |       |---- stats.csv
|   |       |---- hits.csv
|   |       |---- AG24_13.csv
|   |       |---- ...
|   |       `---- correlation_{condition,control}.png
|   `---- z_score/
|       |---- AG24_13/
|       |   |---- enrichment_z_score.R
|       |   |---- stats.csv
|       |   `---- hits.csv
|       |---- AG24_14/
|       |   |---- enrichment_z_score.R
|       |   |---- stats.csv
|       |   `---- hits.csv
|       `---- AG24_15/
|           |---- enrichment_z_score.R
|           |---- stats.csv
|           `---- hits.csv
\end{lstlisting}

For a typical comparison, DELT-Hit yields analysis-ready count tables, generated analysis scripts, replicate-correlation plots, and ranked hit lists that can be inspected before interpretation. In well-behaved examples, enriched compounds are separated from the background by a clear dynamic range, whereas unstable controls or inconsistent replicate structure indicate that troubleshooting is needed before biological conclusions are drawn. Troubleshooting advice can be found in Table~\ref{tab:troubleshooting}.

\subsection{Phase 6: Full-library enumeration and downstream representations}
Optional full-library processing produces files suitable for cheminformatics and machine-learning workflows:

\begin{lstlisting}
campaign/
|---- library/
|   |---- library.parquet
|   `---- properties/
|       `---- library/
|           |---- properties.parquet
|           `---- prop_*.png
`---- representations/
    |---- morgan.npz
    `---- bert.npz
\end{lstlisting}

These outputs support downstream analysis of the complete library rather than only the selection-focused subset. In practice, the expected results are (i) a full enumerated compound table with calculated descriptors and (ii) standardized numerical representations that can be used directly for similarity search, clustering, visualization, or machine-learning models built with common Python frameworks.
